\scanpage{90}
\mspage{89}

\uline{Définition} Soit $X\in\unsure{M|\mathcal M}$. Posons
\[
X^\circ=\left\{Y\in\operatorname{Omb}(X)\ \middle|\ \begin{aligned}
&Y\parallel\partial X\quad\text{i.e. }Y_{\partial X}=\varnothing\quad\text{i.e.}\\
&\deleted{\unsure{\exists Z}}\widetilde Y\cap\operatorname{Omb}(\partial X)=\varnothing\\
&\text{i.e. }\unsure{\exists|\nexists}Y'\leq Y,\ Y'\ll X'<X
\end{aligned}\right\}
\]
\margin{ou $\operatorname{Omb}^{\unsure{\circ\circ|\infty}}(X)\subset\operatorname{Omb}^{\circ}(X)$}
% slanted
\margin{partie fermée de $M,\ll$}

Donc pour $X,Y\in M$, la relation
\[
X^\circ\parallel Y^\circ
\]
\unsure{est} définie, elle signifie
\[
\forall X'\in X^\circ,Y'\in Y^\circ,\quad\text{on a}\quad X'\parallel Y'.
\]
\margin{Attention, \uline{\unsure{je ne|ne me}} \uline{dis pas} que $X^\circ\cap Y^\circ=\varnothing\Rightarrow X^\circ\parallel Y^\circ$, cela \unsure{impliquerait} pour $x,y\in\unsure{L|\ell}$ $x\ne y\Rightarrow x\parallel y$ !}
% slanted
\margin{\uline{NB} $Y\in X^\circ$ signifie \unsure{aussi que} $Y\ll X$, et l'application $\unsure{\varphi|\phi}:\widetilde Y\to\widetilde X$ est constante de valeur le \deleted{\illegible{1}} plus grand élément de $X$. Au lieu de $Y\in X^\circ$, on écrit aussi $Y\in\operatorname{Omb}^{\circ\circ}(X)$, ou $Y\overset{\circ\circ}{\ll}X$ (\unsure{raffinement} \uline{intérieur} de $X$, plus fort que \unsure{raffinement} \uline{immédiat} $Y\overset{\circ}{\ll}X$ i.e. $Y\in\operatorname{Omb}^{\circ}(X)$).}

On voudrait prouver ceci

\inserted{(OK si on suppose At L2, L3)}

\uline{Lemme 1} (?) Soient $X,Y\in M$. Alors
\[
X\parallel Y\Longleftrightarrow\forall X'\in\widetilde X,Y'\in\widetilde Y,\quad\text{on a}\quad X'^\circ\parallel Y'^\circ.
\]

\uline{Dém.} $\Rightarrow$ est clair, car $X'^\circ\subset\operatorname{Omb}(X')\subset\operatorname{Omb}(X)$ et de $=$ $Y'^\circ\unsure{\subset|\in}\operatorname{Omb}(Y')\subset\operatorname{Omb}(Y)$, et on applique At D$_0$.

L'implication $\Leftarrow$ est moins claire, \deleted{Elle} l'est \deleted{\unsure{mais est loin alors à prouver ceci|mais est loin alors d'être claire}} car on doit \unsure{en admettre} L2, L3.
% slanted
\margin{\uline{Cor} $X\parallel Y\Longleftrightarrow\forall Y'\in\widetilde Y,\ X\parallel Y'^\circ$}

\uline{Lemme 2} \inserted{(At \textcircled{\small L1}, L2, L3)} \deleted{Supposons} Soit $Y\ll X$. Alors
\[
\text{a)}\quadY^\circ\overset{\circ}{\ll}X\quad\text{sss}\quad Y^\circ\subset X^\circ.\quad\text{On a}
\]
\[
\text{b)}\quad Y\overset{\circ\circ}{\ll}X\quad\inserted{\text{i.e. }Y\in X^\circ}\quad\text{sss}\quad\forall Y'\in\widetilde Y,\quad Y'^\circ\subset X^\circ.
\]

\uline{Dém.} a) Supposons $Y\overset{\circ}{\ll}X$, prouvons $Y^\circ\subset X^\circ$. Soit $Y'\in Y^\circ$, i.e. $Y'\ll Y$, $Y'\parallel\partial Y$, donc $Y'\ll X$, il faut prouver $Y'\parallel\partial X$ ou encore $Y'_{\partial X}=\varnothing$
